\documentclass[12pt]{report}

\usepackage[a4paper, top=1in,bottom=1in,right=0.8in,left=1.5in]{geometry} % [width=150mm,top=25mm,bottom=25mm]

\usepackage[utf8x]{inputenc}
\usepackage{setspace}
\usepackage{mathptmx} % for Times font
\usepackage{acronym}
%\usepackage[square, numbers, comma, sort&compress]{natbib}
%\usepackage{apacite}
\usepackage[backend=biber, style=apa,citestyle=apa]{biblatex}
\addbibresource{ref.bib}
\usepackage{enumerate}
\usepackage{hyperref}
\usepackage[english]{babel}
\usepackage{amsthm}
%\usepackage{layout} to visually display layout of my page use \layout command

% chapter format
\usepackage{titlesec}
\titleformat{\chapter}[display]{\normalfont\Large\bfseries\centering}
{\chaptertitlename\ \thechapter}{10pt}{}[]

% dummy text
\usepackage{lipsum}

% image related packages
\usepackage{graphicx}
\graphicspath{{images/}} % set image path
\usepackage{subcaption}
\usepackage[export]{adjustbox}
\usepackage{wrapfig}
\usepackage{comment}

% glossaries
\usepackage[acronym]{glossaries}

%codes
\usepackage{listings}
\usepackage{minted}

\newtheorem{theorem}{Theorem}[section]
\newtheorem{corollary}{Corollary}[section]

\theoremstyle{definition}
\newtheorem{definition}{Definition}[section]

\theoremstyle{definition}
\newtheorem{example}{Example}[section]

\theoremstyle{remark}
\newtheorem{remark}{Remark}[section]

\theoremstyle{definition}
\newtheorem{problem}{Problem}[section]

\makeglossaries
\loadglsentries{abbreviations}
\begin{document}
\singlespacing

\pagenumbering{roman}

\input{titlepage}

\chapter*{\Large DECLARATION OF ORIGINALITY}
% \setcounter{page}{2}
\input{declaration}
\addcontentsline{toc}{chapter}{DECLARATION OF ORIGINALITY}

\chapter*{\Large THESIS EVALUATION}
% \setcounter{page}{2}
\input{evaluationPage}
\addcontentsline{toc}{chapter}{THESIS EVALUATION}

\chapter*{\Large PERMISSION}
% \setcounter{page}{2}
\input{permission}
\addcontentsline{toc}{chapter}{PERMISSION}

\onehalfspacing
% \renewcommand{\baselinestretch}{1.5} 

\chapter*{ACKNOWLEDGEMENTS}
% \setcounter{page}{2}
\input{acknowledgements}
\addcontentsline{toc}{chapter}{ACKNOWLEDGEMENTS}

\chapter*{ABSTRACT}
\input{abstract}
\addcontentsline{toc}{chapter}{ABSTRACT}

\tableofcontents
\addcontentsline{toc}{chapter}{CONTENTS}

\listoffigures
\addcontentsline{toc}{chapter}{LIST OF FIGURES}

\listoftables
\addcontentsline{toc}{chapter}{LIST OF TABLES}

\printglossary[type=\acronymtype, title=Abbreviations]
\addcontentsline{toc}{chapter}{ABBREVIATIONS}

%\chapter*{Abbreviations}
%\input{abbreviations}
%\addcontentsline{toc}{chapter}{Abbreviations}

\clearpage
\pagenumbering{arabic}

\chapter{INTRODUCTION}
\label{introduction}
\input{chapters/ch1Introduction}

\chapter{LITERATURE REVIEW}
\label{litreview}
\input{chapters/ch2LiteratureReview}

\chapter{METHODOLOGY}
\label{meth}
\input{chapters/ch4Methodology}

\chapter{DATASETS}
\label{datasets}
\input{chapters/ch5Datasets}

\chapter{EXPERIMENTS}
\label{experiments}
\input{chapters/ch6Experiments}

\chapter{EXPERIMENTAL RESULTS AND DISCUSSION}
\label{results}
\input{chapters/ch7ResultsAndDicussion}

\chapter{CONCLUSION}
\label{conclusion}
\input{chapters/ch8Conclusion}

\appendix
\chapter{HARDWARE AND SOFTWARE CONFIGURATIONS}
\label{appendixA}
\input{chapters/appendixA}


\chapter{MODEL DETAILS}
\label{appendixB}
\input{chapters/appendixB}

%\bibliographystyle{ieeetr}
%\bibliographystyle{apacite}
%\bibliography{ref}
\printbibliography
\addcontentsline{toc}{chapter}{BIBLIOGRAPHY}

\end{document}
